\subsection{Maldición de la Dimensionalidad}
La maldición de la dimensionalidad o \textit{curse of dimensionality}
es un fenómeno que ocurre cuando el número dimensiones de un conjunto de datos crece,
específicamente cuando el muestreo no aumenta lo suficiente. Obsérvese por ejemplo que el
número de elementos en una malla uniforme de un hipercubo requiere un crecimiento exponencial del
número de muestras con el número de dimensiones.

\begin{observacion}
	Dada la ecuación del ATE:
	$$
		\widehat{ATE} = \sum_x P[X = x] \left(P[Y = 1 | X = x, A = 1] - P[Y = 1 | X = x, A = 0]\right)
	$$

\end{observacion}
